% Part: many-valued-logic
% Chapter: three-valued-logics
% Section: goedel

\documentclass[../../../include/open-logic-section]{subfiles}

\begin{document}

\olfileid{mvl}{thr}{god}

\olsection{Logika-logika G\"odel}

Kurt G\"odel ngenalake rerangkening logika $n$ nilai kang
saben logikane ngemot kabeh !!{formula}s kang sah ing logika
intuisionistik lan kalebu ing logika klasik. Iki kang kapisan
narik kawigaten:

\begin{defn}\ollabel{defn:goedel}
\emph{Logika G\"odel $3$ nilai}~$\LogGod$ ditetepake nganggo matriks iki:
\begin{enumerate}
  \item Basa proposisional baku $\Lang L_0$ kang ngemot
  $\lfalse$, $\lnot$, $\land$, $\lor$, $\lif$.
  \item Himpunan nilai kabeneran $V = \{\True, \Undef, \False\}$.
  \item $\True$ mung siji-sijine nilai pinilih, yaiku $V^+ = \{\True\}$.
  \item Kanggo $\lfalse$, kita nduweni $\tf{\lfalse} = \False$.
  Fungsi kabeneran kanggo konektif liyane diwenehake dening
  tabel iki:
  \begin{center}
    \begin{tabular}{c|c} 
      $\tf{\lnot}[\LogGod]$ & \\ 
      \hline  
      $\True$ & $\False$ \\ 
      $\Undef$ & $\False$ \\
      $\False$ & $\True$ 
    \end{tabular}
    \quad
    \begin{tabular}{c|ccc} 
      $\tf{\land}[\LogGod]$ & $\True$ & $\Undef$ & $\False$ \\ 
      \hline 
      $\True$ & $\True$ & $\Undef$ & $\False$ \\ 
      $\Undef$ & $\Undef$ & $\Undef$ & $\False$\\ 
      $\False$ & $\False$ & $\False$ & $\False$ 
    \end{tabular}
    \\[2ex]
    \begin{tabular}{c|ccc} 
      $\tf{\lor}[\LogGod]$ & $\True$ & $\Undef$ & $\False$ \\ 
      \hline 
      $\True$ & $\True$ & $\True$ & $\True$ \\ 
      $\Undef$ & $\True$ & $\Undef$ & $\Undef$ \\
      $\False$ & $\True$ & $\Undef$ & $\False$ 
    \end{tabular}
    \quad
    \begin{tabular}{c|ccc} 
      $\tf{\lif}[\LogGod]$ & $\True$ & $\Undef$ & $\False$ \\ 
      \hline 
      $\True$ & $\True$ & $\Undef$ & $\False$ \\ 
      $\Undef$ & $\True$ & $\True$ & $\False$  \\ 
      $\False$ & $\True$ & $\True$ & $\True$ 
    \end{tabular}
  \end{center} 
\end{enumerate}
\end{defn}

Tabel kabeneran kanggo $\land$ lan~$\lor$ padha karo
ing logika \L ukasiewicz lan logika Kleene kang kuwat, nanging
tabel kanggo $\lnot$ lan~$\lif$ beda ing saben logika. Ing logika
G\"odel, $\tf{\lnot}(\Undef) = \False$. Beda karo logika
\L ukasiewicz lan Kleene, $\tf{\lif}(\Undef, \False) = \False$;
beda karo logika Kleene, nanging padha karo logika \L ukasiewicz,
$\tf{\lif}(\Undef, \Undef) = \True$.

Kaya kang wis diisyaratake dening sesambungane karo logika
intuisionistik, $\LogGod[3]$ cedhak karo logika intuisionistik.
Kabeh proposisi kang sah ing logika intuisionistik dadi tautologi
ing~$\LogGod[3]$; akeh tautologi klasik kang ora sah ing logika
intuisionistik uga dudu tautologi ing~$\LogGod[3]$. Umpamane,
formula ing ngisor iki dudu tautologi:
\begin{align*}
  & p \lor \lnot p && (p \lif q) \lif (\lnot p \lor q) \\
  & \lnot\lnot p \lif p && \lnot(\lnot p \land \lnot q) \lif (p \lor q) \\
  & ((p \lif q) \lif p) \lif p && \lnot(p \lif q) \lif (p \land \lnot q)
\end{align*}
Nanging ora saben tautologi ing $\LogGod[3]$ sah ing logika
intuisionistik, umpamane $\lnot\lnot p \lor \lnot p$ utawa $(p
\lif q) \lor (q \lif p)$.

\begin{prob}
  Wenehana tabel kabeneran kanggo nuduhake yen formula ing ngisor
  iki tautologi ing~$\LogGod[3]$:
  \begin{align*}
    & \lnot\lnot p \lor \lnot p\\
    & (p \lif q) \lor (q \lif p) \\
    & \lnot(p \land q) \lif (\lnot p \lor \lnot q) \\
    & (p \lif q) \lor (q \lif r) \lor (r \lif s)
  \end{align*}
\end{prob}

\begin{prob}
  Wenehana tabel kabeneran kanggo nuduhake yen formula ing ngisor
  iki dudu tautologi ing~$\LogGod[3]$
  \begin{align*}
    & (p \lif q) \lif (\lnot p \lor q) \\
    & \lnot(\lnot p \land \lnot q) \lif (p \lor q) \\
    & ((p \lif q) \lif p) \lif p \\
    & \lnot(p \lif q) \lif (p \land \lnot q)
  \end{align*}
\end{prob}

\begin{prob}
  Sesambungan endi ing ngisor iki kang lumaku ing logika G\"odel?
  Wenehana tabel kabeneran kanggo saben sesambungan.
  \begin{enumerate}
    \item $p, p \lif q \Entails q$
    \item $p \lor q, \lnot p \Entails q$
    \item $p \land q \Entails p$
    \item $p \Entails p \land p$
    \item $p \Entails p \lor q$
  \end{enumerate}
\end{prob}

\end{document}
